% !TeX root = ../main.tex

\section{Construction and pitfalls of a dielectrophoretic trapping setup with fluorescence and iScat modes}
\label{sec:TroubleShooting}

\subsection*{Abstract}
Main message: the key sources of problems and how to circumvent them. 

\subsection*{Introduction}
\label{sec:intro_trouble}
	%	Situation: We want to see protein conformational changes\\
	%	Complication: The changes are small and fast, so we cannot use gold standard: \gls{fl}\\
	%	Question: how can we still detect this? \\
	%	Answer: We are going for \gls{DEP} which is potentially sensitive to this\\
	%	Visualisation: not yet\\
	The Goal of our project is to detect protein dynamics. As mentioned in \hyperref[cha:introdcuintroduction]{the first chapter of this book}, single-molecule detection is the best way to do this. Localising and identifying single proteins and their interactions are mature research fields that emerged around the beginning of the 21st century.  Therefore, numerous single-molecule techniques for studying proteins already exist and have improved significantly in recent decades \cite{joo2008advances,NobelPrizePressRelease2025}. Most of them are \fl -based, like simple epi-fluorescence and can reach super-resolution levels as in \gls{palm}\cite{hess2006ultra,betzig2006imaging} and \gls{storm}\cite{rust2006sub}. When looking at protein dynamics, such as conformational changes and interactions with other molecules, label-based techniques do exist, such as \gls{fret} \cite{jares2003fret,roy2008practical} and optical tweezers \cite{ashkin1986observation,capitanio2013interrogating}, but they all rely on labelling the molecules of interest \cite{huang2009super,weiss1999fluorescence}. \\ 
    Fluorophores have several disadvantages: they can blink and bleach, making long-term measurements difficult, they must be attached to the molecule of interest, which complicates sample preparation and alters the energy landscape of the labelled molecule \cite{yin2015does}. This could affect the observed dynamics if the right part of the molecule were labelled at all and the change were large enough to be detectable. Additionally, the fluorescence lifetime of an excited molecule -- the time the molecule is in the excited state before it goes back to the ground state while releasing energy in the form of fluorescent light -- is $\~10^-8~S$. Anything that happens faster cannot be resolved, which means small domain fluctuations, side chain rotations, small backbone fluctuations and local loop motions\cite{wolff2023mapping,nam2023protein,khodadadi2015protein,mcdonald2013colocalization}. These rapid motions, in turn, influence large-scale conformational changes, ligand binding, allosteric transitions, and folding dynamics \cite{haran2023perspective}. \\ 
    Therefore, label-free techniques like nano-pores\cite{schmid2021nanopores} and \gls{iscat}\cite{hsieh2014tracking,ortega2012interferometric} have been developed. They are able to observe proteins in a more natural state, in vivo, even \cite{Mazaheri24}. However, they are not specific enough to observe single molecules in a mixture, requiring \fl labelling after all, or only giving limited information about the molecules. Additionally, some of the aforementioned rapid motions are only resolved in silica \cite{nicolai2020Nanopore, yusko2016Realtime}. There is a gap in technologies capable of observing single, unlabelled molecules and their dynamics across all time and size scales of change. \\	
	All the aforementioned techniques rely mostly on properties that are not directly linked to the protein conformation, like the location of the fluorophores attached to the molecule and the scattering cross-section. An inherent property of molecules that is directly linked to their structure is the dipole-moment \cite{sarkar2017comparison,piecek2012orthoconic,pereira2018machine,mcgrath2007spatial}
	. It is determined by the folding of the protein chain and the distribution of charged and polar side-chains and therefore changes with the confirmation of the protein and in the face of interactions with other proteins or charged ligands like ions. Research into this parameter is still in its crib \cite{pethig2022Protein}, but techniques using it more or less unknowingly are widespread. \gls{DEP} is commonly used for sorting cells and proteins, in on-chip mixers, etc. The dipole-moment interacts with the inhomogeneous electric field, so it experiences the \gls{DEP}-force (formula in the previous chapter). Measuring this force enables the detection of the dipole moment, which indicates the structure of the protein.

	%	Situation: \gls{DEP} force theoretically can report on protein changes \\
	%	Complication: \gls{DEP} is a small signal and difficult to read out\\
	%	Question: how can we measure the \gls{DEP} force? \\
	%	Answer: Optically using \gls{iscat} (explain what is \gls{iscat}, why it needs shot noise limit, why we still use \gls{fl} and how \gls{DEP} setup works), what parameters do we need for single molecule measurements, why nano gaps  \\
	%	Visualisation: Figure showing schematic of our setup and what results to expect in \gls{proteinDEP} \\
	\begin{figure}
		\centering
		\includegraphics[width=\textwidth]{pictures/trubelshooting/Setup/Setup_v3_1.png}
		\caption{\textbf{The goal.}
			\textbf{A} Schematic of the final setup, with Illumination (laser, optical isolator, telescope, polarisation control, beamsplitter), object table (objective, object table, sample) and detection (filters and pathways for fluorescence (\gls{spcm}) and i\gls{iscat}t (\gls{femto}) detection). 
			\textbf{B} Schematic of \gls{ps} particle detection during trapping process, with object table and zoom in to side view of electrodes and detection volume.
		}
		\label{fig:Setup}
	\end{figure}
	For single molecules this force is too small to read out directly, but we think it could indirectly be quantified through measuring the \gls{DEP}-force by combining \gls{DEP} with a label-free optical technique. 
	We chose \gls{iscat} for this. It is label-free and capable of detecting single protein molecules by measuring the interference of light that they scatter with ther reflection from the glass-water interface of the cover slip and the medium. The scattering cross section depends on the molecular mass of the protein. For a more thorough explanation of this technique see the previous chapter. \\ 
    Herein, we build a combined \gls{DEP} and \gls{iscat} setup, shown in \cref{fig:Setup}. The biggest part is the confocal \gls{iscat} microscope shown in \cref{fig:Setup} A. A $633~nm$ HeNe-laser hits a 90:10 beam splitter, where 10\% of this light get thrown into the high NA objective ($NA=1,4$), which focusses the beam onto the sample. The signal returns through the objective and 90\% of it through the beam splitter onto a dichroic mirror. The scattering signal is focussed through a $40~\mu m$ pinhole to filter out background scattering and finally focussed on a \gls{femto}. \gls{iscat} is complementary to \fl and eases the setup construction immensely (more on that in \cref{alignment}), therefore we want to simultaneously record a \fl signal. The dichroic mirror transmits the redshifted \fl signal which gets focussed onto a fibre acting like a confocal pinhole ($125~\mu m$ cladding diameter) and guides the light to a \gls{spcm}. \\ The confocal signal comes from an object experiencing \gls{DEP} force and entering the detection volume as shown in \cref{fig:Setup} B. The force is produced by the electric field of a pair of triangle-shaped, gold electrodes with gap sizes down to $20~nm$. The triangular shape creates a highly localized electric field between the tips of the electrodes. The nano-scale gap size is necessary to achieve high field strengths ($10^8~\frac{V}{m}$) at low voltages ($1-5~V$). In turn, low voltages are necessary to protect the setup from heating effects by the electric fields, which could cause unwanted flow or denaturation of proteins. By trapping single molecules inside the nano-electrodes and detecting them in our confocal setup we hope to measure the \gls{DEP} force optically and de-convolute it to determine the protein dipole moment (details on this in \cref{chapter:introduction}). 
   
	%    Situation: We want to build a \gls{proteinDEP} setup\\
	%    Complication: this does not work like a regular \gls{iscat} (this needs some polishing)\\
	%    Question: What are complicating factors in our setup? Why are we so sensitive?\\
	%    Answer: long trapping times and shiny electrodes\\
	%    Visualisation: Not needed\\  
	\begin{figure}
		\centering
		\includegraphics[width=0.9\textwidth]{pictures/trubelshooting/Experiments/Experiments_v3}
		\caption{\textbf{Ideal Tests.}
			\textbf{A} Confocal fluorescence scanning images of PSF of gold nanorods in xy, xz and yz direction. Position of the profile plots marked in yellow.
			\textbf{B} Profile plots of the fluorescence xy-scan in x and y (yellow) direction with Gauss-Fit (green). \gls{fwhm}$_x = 331,81 nm$ and $R^2=0,927$ and \gls{fwhm}$_y = 264,37 nm$ and $R^2=0,954$
			\textbf{C} Confocal scattering scanning images of PSF of gold nanorods in xy, xz and yz direction. Position of the profile plots marked in yellow.
			\textbf{D} Profile plots of the scattering xy-scan in x and y (yellow) direction with Gauss-Fit (green). \gls{fwhm}$_x = 259,85 nm$ and $R^2=0,962$ and \gls{fwhm}$_y = 315,1 nm$ and $R^2=0,998$
			\textbf{E} Time trace of gold nanorod (gold) fluorescence overlay over glass background (blue). Filtered with Butterworth at $2000~Hz$ and $10^6$.
			\textbf{G} Time trace of gold nanorod (gold) scattering overlay over glass background (blue). Filtered with Butterworth at $2000~Hz$ and $10^6$.
			\textbf{F and H} Power spectral density of time traces in E and G. It shows the frequencies present in the time trace. Calculated using Welch method.\\
			Ca $140~\mu mm$ thick glass, laser power $=15~\mu W$, \gls{femto} setting $=10^7$
		}
		\label{fig:Experiments}
		
	\end{figure}
	The combination of confocal \gls{iscat}, \fl and \gls{DEP} is novel, therefore, no microscope, fit for our electrodes is on the market. The whole setup is made in house. This is usual for \gls{iscat} setups, which can be made by retrofitting a confocal microscope \cite{jacobsen2006Interferometric}, but more often than not are build in house \cite{ortega2016interferometric}. But we are not creating a regular \gls{iscat} microscope. By trapping our specimens over up to multiple seconds, any noise is integrated in our signal, a problem regular \gls{iscat} users do not have. Additionally, our \gls{fov} contains or is close to the gold nano-electrodes. Any vibrational or acoustic noise is amplified by the moment of these -- in these scales -- huge mirrors, which could over-power our signal. As a result, a lot of effort has been made to mitigate any noise sources, to make \gls{iscat} detection of single molecules even possible.
 
	%    Situation: We want to build a \gls{proteinDEP} setup\\
	%    Complication: Long trapping times and shiny electrodes\\
	%    Questions: What checks do we need to use to confirm a stable setup? \\
	%    ANswer: Scans, \gls{psd}, collumation tools, blocking mirrors, polarimeter, HCl, etc\\
	%    Visualisation: Scan \gls{gnr}, \gls{psd}, diffusion times?, SEM of gap before and after use
    
    To achieve the required resolution for \gls{iscat} single-molecule detection and a low enough background and noise level to mitigate electrode effects and enable measurement of the \gls{DEP} force, we tested the setup using \gls{gnr}. They are an ideal sample as they can be used for both \fl and \scat signalling. To be precise, they showcase luminescence, a plasmonic effect based in the material and shape which also results in a red-shifted signal in addition to their scattering signal\cite{gaiduk2011correlated}. This signal is stable over time, unless the input power is so high that it thermally deforms the \gls{gnr}. By imaging them in \gls{hcl}, they stay stuck to the surface of the cover slip and do not lift off like they tend to do in water. Because they are sub-diffraction particles, their \gls{psf} can be used to determine the alignment quality of the setup in scanning mode \cite{cole2011measuring}. An ideal scan of a \gls{gnr} in both \scat and \fl can be seen in \cref{fig:Experiments} A and C, with the corresponding profile plots showcasing the \gls{fwhm} in x and y-direction. For a particle of this size (SIZE XXX) at our wavelength of $633~nm$ the expected \gls{psf} should have a \gls{fwhm} in x and y of about
    
    \begin{equation}
    	\Delta x_{FWHM}~=~0,61*\frac{\lambda}{NA}~=~275,81~nm
    \end{equation}
    
	According to the Rayleigh criterion. This value is close to our measured \gls{fwhm} of \gls{fwhm}$_x = 331,81 nm$ and \gls{fwhm}$_y = 264,37 nm$ in \fl and \gls{fwhm}$_x = 259,85 nm$ and \gls{fwhm}$_y = 315,1 nm$ in \scat in our ideal scan shown in \cref{fig:Experiments}. The ratio between these values indicates the circularity. A perfect circle has a circularity of 1, ours are lower, but workable ($C_{fl}~=~0,797$ and  $C_{fl}~=~0,82$).\\ \cite{cole2011measuring} have a detailed description on other methods to procure the \gls{psf} and how to interpret it. To identify noise sources in our signal, we parked the laser on a \gls{gnr} and recorded the signal in a time trace. The time traces are always filtered using a Butterworth filter at $2000~Hz$ cut-off frequency with $10^6$ intensity. The filtered traces can be compared with the background of the glass-water interface, or Fourier analysed using the Welch method to extract the \gls{psd}. This shows all the frequencies present during the recording and helps identify specific noise sources based on the frequency range. Mechanical noise, through vibrations and acoustics, will have lower frequencies; the laser or disturbances in the electrical supply will have higher frequencies. \\ The methods used to get to the shown scans and traces can be simple, like blocking or re-routing the light to exclude certain sections of the setup that might cause error, or highly complex, like a polarimeter. Apertures found regular use for alignment, as did a home-brew tube system with semitransparent apertures, for the same reason. Collimation was tested both in the old school manner of measuring the size of the beam across a long distance (meters), as well as using a commercial collimation tool. All of these tools and methods were used in an iterative process of measuring a signal, determining noise sources, improving the setup and repeat. In the end, we created a setup that is \gls{shotnoise} limited and able to detect single molecules.
	
	Lots of expertise was gathered during this process, but the author observed a lack of resources for conveying this kind of expertise to others.
	The goal of this chapter is to give an overview of how our setup is constructed, calibration and optimisation of the optical signal, electrical and mechanical stability of the system and some tools and information on how to construct it yourself using various noise sources as examples.
%\FloatBarrier
\newpage
\subsection*{Process}
	The construction of our setup was an iterative two-step process: Step one entails pacing and aligning all optical components, including the object table and the cage. Step two focuses on identifying and mitigating noise sources. These interventions usually require realignment, which then uncovers other noise sources that were too low before, but now dominate and so on and so forth. This chapter will discuss various noise sources and how to circumvent them. I provide external resources when they exist, and my own experiences and solutions. %Figure \cref{fig:noiseExamples} showcases the major noise sources touched upon in this chapter. 
	
	\subsubsection{Alignment}
    \label{sec:alignment}
		\begin{wrapfigure}{l}{0.42\textwidth}  
			%\vspace{-1em}                      
			\centering
			\includegraphics[width=0.40\textwidth]{pictures/trubelshooting/Missalignment/missalignment_v1}
			\caption{\textbf{Follow the protocol.} 
				misalignment can lead to low or bad signal. See here: x-y-scan of \gls{gnr} in scattering aligned wrong, because scattering was done first, ALSO SHOW XZ OF THIS.
				Laser input power XXX W, FEMTO sensitivity $10^6$.
			}
			\label{fig:alignment}
		\end{wrapfigure}
		The obvious first steps in setting up an \gls{iscat} are building up the air compression table in a suitable room and placing the optical components on it. Many papers, protocols and videos exist on this topic. I want to highlight \cite[McClelland and Mankin]{mcclelland2018optical}, \cite[Jonkman et al.]{jonkman2020tutorial}, \cite[Gemeinhardt et al.]{gemeinhardt2018iscat} and \cite[Cole et al.]{cole2011measuring} (probably others once I have done more reading) here who describe this well and practically. 
        \cite[McClelland and Mankin]{mcclelland2018optical} give an excellent general introduction for any newcomer. They explain all the optical components, scientific principles and various optical setups and their uses. Especially, Chapter 5 gives a comprehensive overview of tips on how to handle an optical setup and how to do it safely, and how to design and construct an optical setup. Supplementing their descriptions with videos is helpfull, the journal JoVE has several helpful tutorials\cite{nichols2011video, felcher2024design} as does ThorLabs \cite{ThorlabsPhotonicsLabVideos2026}\cite{Thorlabs2025OpticsHandling}.
		\cite[Gemeinhardt et al.]{gemeinhardt2018iscat} as a very accessible video on how to construct wide-field transmission(?) \gls{iscat}.
		For general users of confocal fluorescence imaging \cite[Jonkman et al.]{jonkman2020tutorial} has a detailed overview from hypothesis to statistics. It gives tips on which technique to choose, how to optimise a setup to one's needs and a lovely troubleshooting guide for blurry or dim images. Many of their steps apply here as well, even though this research does not deal with confocal imaging further than testing the setup. And finally \cite[Cole et al.]{cole2011measuring} has a detailed description on how to generate and interpret a \gls{psf}.
		%	Situation: We need to illuminate our sample
		%	Complication: The illumination needs to be straight without coma and with circular polarization
		%	Question: How can we ensure a perfectly aligned beam and circular polarization?
		%	Answer: %hitting all mirrors in the centre, %surface coating, %beam splitter and %polarization and angles, %chasing evenness
		The above-mentioned resources are extensive and helpful, but I still want to mention a few learnings I gained through this process, which are details usually not mentioned anywhere else. I will roughly walk through the illumination pathway from the laser to the sample and describe a few interceptions to note to make alignment exact and easy. %The beam needs to stay at the same height from the laser to the last mirror into the objective and hit all optical components as centred as possible. This maintains polarisation and avoids aberration -- especially \gls{coma}. Lenses and pinholes need to be co-linear with the beam path, and mirrors should be at a $45\deg$ angle, more details on how to accive this can be found in \cite[McClelland and Mankin]{mcclelland2018optical}. 
        %This is achieved easiest if you choose a straight path along the screw holes of the optical table and place everything centred on this line\cite{mcclelland2018optical}. Careful: Mirror mounts stand out from their poll, so to be above the line, you intend to follow the poll has to be off the line. To do this, always use 2 mirrors to guide the beam in a new direction or a collinear assembly like the beam expander. 
        While \cite[McClelland and Mankin]{mcclelland2018optical} describe in easy words how to plan any optical setup and place the optical components, however they did not mention to pay attention that all mirrors are the same, with the same polarisation-neutral silver coating. Otherwise, the polarisation could change unpredictably after reflection. They also recommend having all mirrors at 45$^\circ$, which we agree with. The exception to this rule is the beam splitter, which can only accept angles less than $10\deg$; otherwise, it is polarisation-active. Therefore, the mirror towards it and the beam splitter require larger and smaller angles with respect to the beam, respectively. Polarisation is relevant because our alignment sample \gls{gnr} and the electrodes in our sample are anisotropic, so they require circularly polarised light to ensure the signal is only size-dependent.\\ 
        The key section to invest the most time in is the section between the beam splitter and the sample. Especially here, everything needs to be aligned as well as possible. The beam splitter needs to be carefully positioned to guide the light to the next mirror, keeping it perfectly straight along the line. This mirror needs to be perfectly aligned with the corner of the line; otherwise, the beam will not hit the next mirror -- the most important one -- centred. This mirror needs to be placed exactly in the centre of the object table; otherwise, you risk underfilling the objective. The beam going up cannot be angled relative to the objective axis, otherwise it will show \gls{coma}. A home-brew tube construction with frosted apertures can aid beam walking\cite{felcher2024design}, but in the end, the sample reflection is the true reporter of \gls{coma}. To ensure the \gls{coma} is not caused by the sample or objective being at an angle, make sure the object table and all that is mounted on it are as parallel towards the microscopy table as possible. This can require filing down sections or adding tissues and rubbers to even them out. The tubes on which the objective is mounted need to be screwed so that the objective is not mounted at an angle when it is fixed. This work is minuscule, but once the object table is perfectly even, you no longer need to adjust it and can align everything else to it. With all of these interventions, you can be sure the illumination is as good as possible and ready for single-molecule detection. 
		%	Situation: All the optical parts are in place and stable, you have good illumination.
		%	Complication: Right Signal still needs to reach the detectors.
		%	Question: How can we get the right signal to the detectors?
		%	Answer: %if illumination is done well, it is easy, but you need careful work, %right order of alignment (first align \gls{fl} then scattering), what does the pinhole do? Why the fiber?, filters\\
		%	Visualisation: Maybe have a figure on alignment here? -> maybe you find an external source for this.
		Once all systems and optical components are in place and aligned to ensure the beam is at the same height and centred on all, especially the objective to avoid coma, the illumination is complete. The same principles described in the sources above hold for the detection path as well, but given our combination of \fl and \scat in reflection mode, I want to describe how we align this section of the setup. All of this work is done on the collimated, reflected light from the sample, as it passes through the beam splitter. As the beam splitter is quite thick and the back side is slightly angled, the beam is displaced and will leave the planned pathway. This can be compensated by two mirrors in the respective detection pathways. Now the alignment towards the two detectors can begin. \\ 
        The key here is that in the end, both channels' signals overlap in space. To test this, we use \gls{gnr} again. First, couple the reflection to the \gls{spcm}. The first step here starts with the \gls{spcm}. We used a fibre-coupled detector, so the fibre acts as our pinhole. Before any light can be directed towards the fibre, it needs to be aligned with its imaging lens. To do this using an optical cage system with the imaging lens attached to the fibre mount is advisable. Attach a fibre-coupled laser to the fibre mount and ensure the transmitted beam does not display \gls{coma} by adjusting the position of the fibre. After the fibre is in the correct position relative to the imaging lens in the cage, fix this position and do not change it for any further alignment. position the imaging lens so the beam is collimated, and the fibre is at the lens's focal distance. This cage system is now internally aligned and can be moved freely on the setup without requiring adjustments. When placing it anywhere, ensure that it is as parallel to the planned beam line as possible. Now, the detection alignment of the sample signal can commence. \\
        Leave the dichroic in place for this, as it also slightly displaces the beam. It might be necessary to increase the laser power to see the laser. Because the \gls{spcm} is very sensitive and can break when overloaded, keep it unplugged for the general alignment. Remove the fluorescence filters for now. Using two mirrors, beam walk to the correct height and centre it into the imaging lens and fibre system. This system should be co-linear with the beam, which can be achieved by back-reflecting a second laser through the fibre and ensuring both beams overlap along the entire beam path between the fibre and the beam splitter. Now the \gls{spcm} can be reconnected -- make sure, when switching it on, that it is not overexposed -- and the beam coupled to it using the two mirrors, maximising the detected signal. It will be necessary to reduce the beam power during this process until it is barely detectable by a power meter. The right mode coupled to the fibre is a sharp peak; if the peak shows local maxima, refocusing might be necessary. Once the coupling is strong enough, use the reflection to find a particle and maximise the signal again. If the particles show up as negative signals in the scanned image, the focal plane is not correct, or the wrong maximum was used to optimise, detune the coupling and try again, until you detect a positive \gls{psf} -- this one does not need to be perfect. Now the \fl filters can be replaced. If there is no \fl signal, remove them again and recouple the scattering. If there is a \fl signal, take note of the \gls{psf} in xy, xz and yz direction. Optimise the position of the mirrors and objective for both maximal signal and ideal \gls{psf} -- see \cite{cole2011measuring} for a better description of the ideal \gls{psf}. If this proves too difficult, the problem might be further upstream. Only once a good \fl \gls{psf} has been achieved, may you continue with the alignment of the \scat detection. \\ 
        The two mirrors adjusting the beam height and direction in \scat are the dichroic and another mirror into the imaging lens. First, align the imaging and the collimation lens with each other. After, align the following mirrors and the collector lens to the \gls{femto}. Now, place the confocal pinhole between the imaging and collimation lenses and maximise its transmission. After everything is in place, optimise the \gls{psf} using the pinhole x,y and z-position. The relevant parameters are: Signal strength, \gls{psf} shape and overlap with the \fl signal. Both channels are now ready for experiments\cite{nichols2011video}. \\ 
        The reason for me to stress this process -- first \fl then \scat detection -- is that for localisation, the \fl signal is more reliable than \scat. The expected signal can be estimated using the quantum yield, it is almost free of background, and only the particles intended to be measured will be detected. Scattering, on the other hand, can be misinterpreted by using the wrong mode, optimising on a dust particle or an Airy ring instead of the actual particle. Therefore, the best way to align a combined setup like ours is to first align using \fl and only then align \scat.
		
		\paragraph{To summarise}
			 The \gls{psf} is the guide by which to judge the quality of your alignment. The key levers are the beam leading into the objective in illumination and the coupling through the fibre and the pinhole in detection. If there are more issues with your signal, find a short troubleshooting guide on this in the appendix.
    
    \subsubsection{On the table (mechanical noise)}
    \label{sec:mechanical}
        \begin{figure}
        	\centering
        	\includegraphics[width=\textwidth]{pictures/trubelshooting/ObjectTable/ObjectTable_v3}
        	\caption{\textbf{Stability of objective and sample plane is key.} 
        		\textbf{A} Schematic of the final version of the objective mount. It has six screws to prevent misalignment from the centre of the object table and ensure maximal vibrational stability. The whole construct of the table and objective mount is centred on the beam, ensuring centred illumination of the objective's back focal plane. This is relevant for minimal focal displacement and even illumination of the sample. To ensure the objective is mounted parallel to the table, rubbers are placed between the object mount and the object table. 
        		\textbf{B} fluorescence time trace and according \gls{psd} showcasing the effect of an \textsl{eliot stage} (x-y-z manipulator). The mount is decoupled from the object table, and the objective is mounted on a long arm. The decoupling causes all vibrational differences between the objective and the sample to show as a big noise pattern in the \gls{psd}. Input power $=~105~\mu W$, \gls{gnr} in gold, glass in blue.
        		\textbf{C} scattering time trace and \gls{psd} on \gls{spcm} showcasing the effect of a home-made objective holder mounted with a custom L-mount to the object table. The objective is now coupled to the table, reducing the vibrational noise. Still, this is not ideal, as the mount is not stable enough and picks up other vibrations. Laser power not detectable
        	}
        	\label{fig:objectTable}
        \end{figure}
%        Situation: We are building a highly sensitive optical device to detect very low signal and signal changes. \\
%        Complication: Any shift of the sample or the laser, will create noise overshadowing our signal. \\
%        Question: Where can this noise come from and how do we negate this? \\
%        Answer: good object stage (%placing lasers and lenses, %locking mirrors and how to test if they are, %sticking down cables, object stage, % pole height considerations)\\
%        Visualisation: the sample holder is an important part of this and will be highlighted in this matter. \\
        As mentioned before, due to the electrodes in our \gls{fov} and the low signal in \scat, the setup is extremely sensitive to mechanical noise. Any movement of optical components that could shift the beam or from the electrodes can produce noise, higher than our signal. The sources for this are numerous, so I will highlight them here along the beam path in the illumination. \\ 
        The Laser is the first optical device in the beam path. It should be mounted stably and even, for good alignment and to avoid drift over time. Our HeNe-laser does not require cooling, but if your laser does, turn the ventilation off for measurements. \\ 
        Next, the beam hits various mirrors, lenses and a pinhole. It is worth spending the money on high-quality mounts with minimal drift and locking mechanisms. Magnetic or flippable mounts are not advisable. These interventions prevent drift and minimise vibrations. All optical components should be mounted on solid pedestal pillar polls and mounted on the table with clamping forks. Adjustable polls could also drift and have more room for vibrations. The beam height should be high enough to allow handy working on the setup and to allow adjustment of the poll height in accordance with the attached mount, but not so high as to necessitate long polls that could act as antennas and pick up environmental vibrations. We opted for a beam height of $8~cm$.\\ 
        A small intervention that can help a lot is sticking down cables on the table using tape. Similar to long polls, can the long cables connecting the cloud to the table capture environmental vibrational noise and transfer it to the components they come into contact with. To test if the noise in your signal is coming from a particular component, you can touch it, tape it, or tighten the locks again, if it has been the source, the noise will reduce in the \gls{psd}.\\\\
        A key component in this is the object table, as mentioned before. The previous section describes how to ensure object table evenness. Additionally, this is the biggest lever for protecting against vibrational noise, as it connects the objective -- the key component for beam guidance -- to the sample -- the key component of the whole setup. I want to take some time to highlight the improvements possible in this assembly, as shown in \cref{fig:objectTable} B. Other \gls{iscat} setups will have an objective mounted fast to the table to avoid vibrations\cite{ortega2016interferometric}. Because our samples have various heights, we need our objective mounted on a z-stage, which does induce some noise. Our Initial design was the object table -- consisting of 4 polls, a breadboard with a manual and a PI-stage mounted on top of it -- and an objective mount -- consisting of an \textit{Eliot} manipulator with a custom Z-shaped extender arm that holds the objective. This is a very unfortunate design because it decouples the objective from the sample. This means that any vibrations from the object and the objective are out of tune with each other, causing the noise to be amplified. Additionally, the extender is another antenna amplifying any unwanted movement, especially in the z-direction. This shows as a high level of low-frequency noise in the \gls{psd} in \cref{fig:objectTable} B. Therefore, we decided to improve this construction by mounting the objective mount on the bottom of the object table breadboard using a differential screw to allow small movements, along with an L-shaped metal piece. This has the advantage that the objective and object are coupled to each other via the breadboard, their vibrations are in tune and do not amplify each other. This becomes clear when looking at the scales in the \gls{psd} in \cref{fig:objectTable} C and D: while in the \textit{Eliot}-configuration the noise level was around $6*10^{-6}~\frac{1}{Hz}$ with maxima at $140~Hz$ and $200~Hz$, but with the first improvement the level drops down to  $4,5*10^{-7}~\frac{1}{Hz}$ with the only relevant low frequency maxima at $200~Hz$. This improvement is already good, but we hoped to do better. The L-shaped piece was not stable enough, and the long differential screw acted like an antenna. The design we came up with is shown in \cref{fig:objectTable} A. Again, the objective mount is attached to the bottom side of the breadboard, this time using 4 screws and 2 L-shaped clamps. Now the mount is stabilised from all four sides, and the differential screw is shorter. The result of this new design is shown in \cref{fig:Experiments} E-H. Extra vibration damping was achieved by mounting the entire construction on a large half-centimetre-thick piece of rubber, with an additional breadboard under the polls. This protects the whole construction from vibrations transmitted into the object table via the optical table. The $200~Hz$ peak persists, but the noise was low enough to detect single particles.
        \paragraph{To summarise} Especially the \scat signal is low and vibration-sensitive due to the relatively large electrodes. Therefore, vibrational stability is key. The mechanical parameters of the setup need to be chosen well. The objective mount is a key component that can be modified for maximum stability. To test if the noise is caused by vibrations, stomping on the floor or tapping on the table or optomechanical piece in question can help determine if this is the cause of the noise.
        
    \subsubsection{The cage (acoustic noise)}
        \label{sec:cage}
    %    	Situation: We have a mechanically stable setup\\
    %    	Complication: There is still noise causing our sample to shift in space. \\
    %    	Question:  Where can this noise come from and how do we negate this? \\
    %    	Answer: %A very simple test (clapping) shows this is acoustic noise (difficult to separate from mechanical noise). We can reduce it by isolating our setup as much from the environment as possible (%cage construction,% notes on the cloud and AirCo and microphones, %floating tables, %isolated bright field). \\
    %    	Visualisation: the cage was the first step to isolating the setup, together with the table, maybe also with the acoustic foam if it does anything.	
    	The setup, as it stands now, is aligned well, with a nice \gls{psf} and is mechanically isolated against vibrations from the setup and the cloud. But even with all components in line and screwed tight and a perfectly even and coupled object table and objective, there will be another form of mechanical noise: acoustic noise. I decided to give this its own section because it is a more complicated problem with different tactics for finding and solving it. On top of that, it is not fully solved for us either at the time of writing.\\  
        The majority of acoustic noise stems from the environment of the lab. Air-conditioning, fridges, walking and talking can all show up in the \gls{psd}. To test if your setup is acoustic sensitive, clap very loudly, if you see a spike on the real-time \gls{psd}, your setup is sensitive. For us, acoustic noise tends to trigger a spike at $100~Hz$, but when playing very loud music, other frequencies can appear as well. Narrowing down which parts of your setup are sensitive is harder, though. A general solution is to cage the system using cardboard. Not only does this create a safer working environment and shield background light, but it also shields some of the acoustic noise present in the lab, as shown in \cref{fig:acoustic}.\begin{wrapfigure}[28]{r}{0.42\textwidth}  
    		\centering
    		\includegraphics[width=0.40\textwidth]{pictures/trubelshooting/Cage/Cage_v4}
    		\caption{\textbf{Acoustic isolation.} 
    			\textbf{A} Picture of the caged setup. The cage protects the setup form background light and reduces the sensitivity to acoustic and vibrational noise.
    			\textbf{B} \gls{psd} of noise in setup measured on \gls{gnr} before the enclosure (light gold) and after (dark gold).
    			Laser input power 15 $\mu W$, FEMTO sensitivity $10^7$
    		}
    		\label{fig:acoustic}
    	\end{wrapfigure} 
        Detaching the overhead rack or cloud from the microscopy table is a potential solution as well \cite{mcclelland2018optical}. The cloud's surface is large, quite heavy, and high up. If your lab allows for a setup to disconnect the two, opt for this. Some groups use marble slabs under their object table to protect from this \cite{gemeinhardt2018iscat}. \cite[Jonkman et al.]{jonkman2020tutorial} also has some more ideas on setup stabilization. \\
    	The final note in this section is about an optical assembly not mentioned before: The bright field. To identify where on the chip you are and to position the gaps close to the laser, a transmission \gls{bright field} pathway is advisable. However, the transmission element of this assembly requires it to hang above the sample. A long pole with heavy optics attached acts like a microphone, transmitting every sound to the table. To avoid this, mount the \gls{bright field} only when needed or hang it from the cloud. 
        \paragraph{To summarise} 
        Acoustic noise is related to mechanical noise, but not quite the same. To detect it is easy, to mitigate it is harder, and it depends on your lab environment. A cage can help eliminate many problems, including the acoustic noise. Any tall constructions should not be attached to the table. To test whether the observed noise is caused by acoustics, playing a loud tone or clapping next to the setup can help identify it.

    \subsubsection{laser noise}
    \label{sec:laser}
        %Situation: All the optics and devices are in place and electrically stable \\  
    	%Complication: Your signal can still be noisy and contain unnecessary background and unideal signal \\
    	%Question: Where can this noise come from and how do we negate this? What can we do to get a perfect signal?\\
    	%Answer: %Single mode fibers can do mode hopping and jump in intensity, %use a non-fibre coupled laser with an optical isolator instead. Make sure the beam shape is Gaussian and polarisation controlled (why not to use single mode fibers, optical isolators, considerations for the laser choice, ideal beam shapes, polarisation)\\
    	%Visualisation: A picture of the laser setup with beam shape (on a paper), time trace of mode hopping and time trace of our setup. (this is not much of a novelty, does maybe something else is more relevant?)\\
        Jonkman et al. \cite{jonkman2020tutorial} recommend carefully selecting the equipment. This holds especially true for the laser source. Even with an acoustically and mechanically stable setup perfectly aligned, the laser can still disturb the measurement. Unstabilized lasers can produce laser noise, which appears as multiple peaks in the \gls{psd} at high frequencies. This is a known and common issue \cite{paschotta2009noise}, so check on this when ordering. \\ 
        But even high-quality stabilised and cooled lasers can throw off the illumination. The first iteration of our setup included a fibre-coupled laser. However, after warming up, the laser would shift mode ever so slightly. The attached single-mode fibre would then not accept the input light, so the light intensity would shift at random intervals. A time trace of how this looks on a detector directly is shown in \cref{fig:laser} B. The variance is huge, but inconsistent, and so is the laser signal itself. The reason for this was most likely back reflections from any surface in the setup, which entered the laser cavity \cite{heumier2003applnote}. After switching to a \textsl{HeNe}-laser and installing an optical isolator to protect the laser from back reflections, our signal has stabilised as can be seen in \cref{fig:laser} C.
        \begin{figure}
    		%\vspace{-1em}                      
    		\centering
    		\includegraphics[width=\textwidth]{pictures/trubelshooting/Laser/Laser_v3}
    		\caption{\textbf{Laser setup.} 
    			\textbf{A} Picture of laser output with optical isolator for stability and telescope for beam shape. Gaussian beam shape is highlighted with an image with Gaussian for of the telescope output.
    			\textbf{B} Time-trace of fibre couple laser without optical isolator displaying mode hopping. SAMPLE IF IT WAS THERE
    			\textbf{C} Ideal time trace after telescope with the setup shown above.}
    		\label{fig:laser}
    	\end{figure}    
    	\noindent
        In fact, most lasers require an optical isolator to protect the laser generation system \cite{heumier2003applnote, mcclelland2018optical}. This optical device consists of two polarisers and a Faraday rotator. That way, the reflected light has the wrong polarisation to pass through the input polariser. Mounting this component at a slight angle directly in front of the laser will protect the laser cavity, increasing the laser lifetime and stability \cite{optisolator}. \\\\
        Confocal microscopy setups work best with a spherical Gaussian beam profile \cite{marvin1961microscopy, carter1972electromagnetic, hauer2016spot}. Gaussian beams are less sensitive to aberrations and do not have side lobes that could broaden the \gls{psf}\cite{sheppard1994influence}, the polarisation is more consistent \cite{carter1972electromagnetic}, and the signal can be focused to a well-defined airy-disc for background elimination \cite{marvin1961microscopy}. This means spherical Gaussian beams can focus to symmetrical, diffraction-limited spots with circular polarisation, thereby maximising background suppression. Most lasers automatically produce a Gaussian beam, but to ensure the beam shape is as symmetric as possible, we installed a beam-cleaning system or beam expander. It consists of an imaging lens (Olympus objective) that focuses on a pinhole (size = 1 Airy disc) and a collector lens that collimates the beam. An aperture blocks the higher-order Airy rings, leaving a symmetrical beam shape that is a perfect spatially limited Gaussian beam. This has multiple advantages, in addition to the beam profile: the beam is now collimated, which allows the optical setup to be modular and easily reconfigurable, and the beam is expanded to the correct size to ensure the objective back aperture is filled, enabling us to use its full \gls{NA}. The polarisation needs a predictable \gls{psf}, as we use it in \cref{sec:alignment}, which is dependent on a high-quality Gaussian beam \cite{sheppard1994influence, hauer2016spot}. \\
        The predictable illumination is especially relevant for our anisotropic samples (\gls{gnr} and nano-electrodes). Especially for the \gls{gnr}, the signal for all particles should be the same. We achieve this by circular-polarisation of our illumination, so that the interaction of the surface plasmon and the illumination is the same for all particles, no matter their orientation \cite{gaiduk2011correlated}. Potentially, we could reduce the background from the electrodes by placing an analyser in the detection pathway and filtering their contribution to the reflection. Further ways of maintaining polarisation were mentioned in \cref{sec:alignment}.
        \paragraph{To summarise} 
        Choosing a high-quality laser source pays off for good signal quality and nice \gls{psf}. Lasers require an optical isolator. Good \gls{psf} require beam cleaning using a beam expander and polarisation control. The laser noise can be identified as a constant sharp peak in the \gls{psd}.
    
    \subsubsection{Electrical noise}
    \label{sec:electrical}
         % Situation: the setup is mechanically and acoustically stable\\
    	% Complication: every now and then signal goes weird, its not acoustic or mechanical\\
    	% Question: Where can this noise come from and how do we negate this?\\
    	% Answer: No matter if input our output, irregular signal fluctuations are often cause by bad electrical input into the device. (switching cables is the first thing, isolated power input)  \\
        % Visualisation: Visualisation is not needed for this as the main take away is "switch the cables". we can refer to the intro.
        This setup uses highly sensitive instruments, and, due to the dual-detection mode, the active laser source, positioning control, and \gls{DEP} application, the technical background outside the optics is quite complex. This complexity introduces more opportunities for failure, which -- in this chapter -- means introducing noise. Several times during this project, we observed irregular signal fluctuations that could not be explained by acoustics, the laser, or setup stability. Upon further investigation, it would turn out that changing the power supply cable could stop this. A good indicator is if the irregularities appear even if the setup has not been used or if a cable is touched. Evidently, even the best machinery can only work with a reliable power supply, so special attention is needed here. \\
        The first step we took to protect our signal from electrical power-line noise was to introduce a power line-filter (BRAND). This filters out frequencies of other appliances connected to the general power supply, and the eigenfrequency of the power supply itself ($50~Hz$). Those frequencies could appear in our \gls{psd}, adding noise in the low-frequency range and potentially interfering with our signal determination. \\
        The second step is an ongoing awareness of all the cables in the setup. Power cables, information transport cables, input, output, etc. Here is a list of advice on preventing noise coming in through the power supply or input/output of your devices: 
            \begin{itemize}
                \item Some of the connections between devices are home-made, and their isolation needs to be checked regularly. 
                \item It is better to have one long information cable than multiple short ones.
                \item If you have to use a connector, make sure it has no open ports. 
                \item If you find a faulty cable, dispose of it immediately. 
                \item Cable management should be tidy to prevent knots that could cause tension on the plug. 
                \item Cable management should minimise cable interactions and prevent accidental physical contact. 
            \end{itemize} 
        In general, our first step in troubleshooting a new noise source was to replace the device's cables. 
        
        \paragraph{To summarise}
        Power supplies and information cables can be a source of noise, as well as the mains electricity. Luckily, this is easily mitigated with the interventions described above, primarily a power line filter and vigilant setup management. The easiest way to check whether the noise is caused by a cable is to replace it. 
        
	\subsubsection{detector background}
	\label{sec:detector}
		\begin{wrapfigure}{r}{0.42\textwidth} 
    			\centering
    			\includegraphics[width=0.48\textwidth]{pictures/trubelshooting/ShotNoise/ShotNoise_v6.png}
    			\caption{\textbf{\gls{shotnoise} limit of \gls{gnr} signal.} 
    				\textbf{A} Laser power into the objective vs detected \scat (yellow) and \fl (orange) signal of %glass-water interface (colour), 
                    \gls{gnr} %and electrode gaps (colour). 
                    Linear fit with $R^2 = 0,949$ and $R^2 = 0,978$ respectively. \gls{femto} sensitivity $10^6$.
    				}
    			\label{fig:ShotNoise}
		      \end{wrapfigure}
		Some noise sources cannot be mitigated. Even the most stable setup, with acoustic isolation, stable lasers and no electrical noise, will have shot noise or detector background. Other background suppression from the room and stray laser scattering can be achieved by boxing the setup, the laser, and its beam expander, and caging the \gls{spcm}. But shot noise is an inherent effect of the particle nature of light, in which, in a stream of photons, the number of detected photons is not constant but Poisson-distributed. For \gls{iscat} this is a very important threshold, because the detected signal is so small compared to its background \cite{ortega2016interferometric}. The background can be subtracted from the signal as it is constant, but any other noise source could overpower this signal, so nothing but background and shot noise should be present behind the detected signal \cite{ortega2016interferometric}\cite{lin2014shot}.\\
        %Situation: You are stably illuminating an isolated setup\\
		%Complication: there will always be background\\
		%Question: Where can this noise come from and how do we negate this? \\
		%Answer: you can fight as much background as possible (darkening room, separating laser, isolating fl sensor, polarisation filters and polarisation of mirrors) but there will always be \gls{shotnoise} (How to protect your detectors, what you cannot get rid of) \\
		%Visualisation: Figure showing we are \gls{shotnoise} limited (laser power vs signal on \gls{gnr}), compare with \gls{femto} and \gls{spcm} specs.  \\
        If a setup is shot-noise-limited, this can be determined by calculating the signal variance at different input laser powers. If the variance scales linearly with the laser power, the signal is dominated by shot noise. Other noise sources may still be present, such as laser and detector noise, but they are negligible compared to shot noise \cite{lin2014shot}. This requires the illumination to be sufficiently bright to deliver sufficient photons \cite{lin2014shot}. \ref{fig:ShotNoise} shows that both our \scat and our \fl detection variance linearly increase with the input power, proving our setup to be shot noise limited. 
        \paragraph{To summarise}
        Both the laser source and the \gls{spcm} should be caged to limit background. Shot noise, however, cannot be mitigated. A good \gls{iscat} is shot-noise-limited. To test whether a setup is shot-noise-limited, a signal at different power levels needs to be recorded and analysed. 
		
    \subsubsection{Sample}
    \label{sec:sample}
        %Situation: You have stable illumination and detection\\
	    %Complication: you need a sample\\
	    %Question: What should you look at and how do you prepare that?\\
	    %Answer: Depending on the experiment we use various sample (mirrors and glass/water for alignment, simple glass for floaters, electrode samples with tween fr \gls{ps} particle trapping, electrodes with anti-fouling for \gls{si} and protein trapping), this is how we prepare a sample and how it can go wrong (leaky wells, touch on anti-fouling, sample maintenance (prevent drying out blow ups etc)), I will only touch on \gls{fab} as this was not part of my thesis.    	% handling of those chips - RCA, protection from the breakdown  = all the protocols we have developed\\
	    %Visualisation: No visualisation needed, this is very basic. we can reference the first figure. \\
        With the setup fully functional, we need to consider which samples to use and for what purpose. As mentioned above, for alignment, a glass or a mirror is required in the sample plane, later \gls{gnr} in \gls{hcl}. The glass used should be as clean as possible. We used sonication in organic solvents and plasma cleaning, sometimes \gls{rca} cleaning was done  as well.\\
        \begin{wrapfigure}{l}{0.42\textwidth}  
    	    	\centering
    	    	\includegraphics[width=0.38\textwidth]{pictures/trubelshooting/GlassThickness/GuiPhase_v2.png}
    	    	\caption{\textbf{Thickness does matter.} 
    	    		\textbf{A+B} Simulation of \gls{psf} in scattering of a spherical particle on \textbf{A} $170~\mu m$ thick glass and \textbf{B} $200~\mu m$ thin glass. 60x Objective, NA 1.4 and a wavelength of $632~nm$ Simulation and design with the open-source software PSFLab7 developed by \citeauthor{nasse2010realistic} Michael J. Nasse, Jörg C. Woehl. \cite{ton2024nano}.
    	    		\textbf{C+D} Detected \gls{psf} of equivalent particle on \textbf{C} $XXX~\mu m$ thick glass and \textbf{D} $XXX~\mu m$ thin glass. 
    	    	}
    	    	\label{fig:GuiPhase}
	          \end{wrapfigure}
        To measure the sample \gls{psf}, \gls{gnr} were applied to a cleaned glass surface in 12\% HCl solution, as described above. \\       
        When determining the dwell-times of freely diffusing particles, a clean glass and a Tween 20 solution would suffice. For floating proteins, the glass requires an antifouling layer, more on that in \ref{sec:surfaceCoatings}. \\
        Trapping requires a more complicated sample preparation. The chip manufacturing was not part of this work and will not be elaborated on. At the end of this multistep clean-room process, a XXX by XXX chip is fabricated on XXX Glass with nine \gls{DEP}-gaps. Due to their nanoscopic size, they are very fragile and can explode from static discharge when handled. To protect them from this, they are connected by three gold short-circuiting rings on the chip's edge. After the chip was finished, it had to be imaged in \gls{sem} to determine the exact gap sizes. This process requires applying a contrast solution, which must be washed off afterwards with water. Doing 30 min of \gls{rca} cleaning or 30 s of \gls{piranha} application is advisable as well. A clean \gls{pdms} well was attached to the chip by plasma activating the well. The placement needs to be as centred as possible, otherwise the chip will not fit into the \gls{pcb}. The activation should only take $30~s$ in plasma; otherwise, the bind is too strong, and removing the well can break the chip. Well can be reused when cleaned with acetone or \gls{rca} If the well or chip surface is dirty, or the well is not attached to the chip immediately upon activation, adhesion might fail. The solution applied to the chip will leak out and compromise the experiment, as the aqueous film will short-circuit the electrodes. If sticky particles, such as silica beads or proteins, are part of the experiment, now is the time to apply antifouling; more on this in \ref{sec:surfaceCoatings}. Next, the chip is mounted into our custom \gls{pcb}, which can control each electrode individually. The \gls{pcb} is mounted on the object table and connected to the rest of the electrical setup, in parallel with a $50~\ohm$ resistor, to protect the gaps from discharge. The \gls{pcb} also encompasses a ground loop, which also protects the gaps from exploding due to discharge. Before starting an experiment, it is essential to ensure that all gaps are properly grounded to prevent accidental damage. Only now can the electrodes be isolated from each other by scratching the protective gold rings using a diamond cutter. The chip is now in a very fragile state and should not be handled if the ground of the \gls{pcb} is not connected to the gaps. Even when adding solutions with a pipette to the well, discharges may cause the electrodes to explode, rendering them unusable. For particle trapping, we used a Tween 20 solution to prevent particle agglomeration or adhesion to surfaces. But even this can fail. If a particle gets stuck in a gap, apply a high-frequency voltage and try agitation via scanning or pipetting. After switching the cut electrodes into active mode and removing the  $50~\ohm$ resistor, trapping experiments can commence. \\
        
      
	    %Situation: You now have a working setup\\
	    %Complication: Your scattering signal is too low\\
	    %Question: Where can this noise come from and how do we negate this? \\
	    %Answer: Your GUI phase needs to be right (I need to read up on this) which is dependent on your glass thickness. \\
	    %Visualisation: simulations showing the change in signal and \gls{psf} based on glass thickness. Show the reality of this using thick and thin \gls{psf} of \gls{gnr}.\\
        The \gls{psf} has been mentioned a few times before in this chapter. As we are building a confocal setup, a good understanding of our focal spot is essential. Because the detected \gls{psf} is a convolution of the illuminated object and the objective's focal spot, when determining the \gls{psf} based on a sample, the sample needs to be taken into account too. Specifically, the substrate thickness. Most objectives are calibrated for a coverslip thickness of 1,5. If there is a thickness mismatch, the focal spot will be shifted, and its shape will change \cite{nasse2010realistic}. When aligning the setup, it is therefore important to use the right coverslip thickness. Our chip fabrication suffered from this problem, too. The chip thickness was inconsistent across batches, leading to suboptimal illumination and signals below the detection limit. We now measure each chip's thickness before delivery to ensure maximal signal efficiency. 
        
        
	    	
        \paragraph{To summarise}
        Nanoelectrodes are sensitive to electrical discharge, so special care is required when using them in experiments. Several custom solutions are used here to assemble the samples. Special attention needs to be taken towards the glass thickness for coverslips and chips. 
        
    \subsubsection{program considerations}
    \label{sec:program}
    	%Situation: The setup with the sample is complete and ready\\
    	%Complication: You need to read your data out\\
    	%Question: What do i need to keep in mind when analysing data generated using our software? \\
    	%Answer: there are physical limitations to what a stage can do (it has to speed up) and there are compromises to be made for the sake of an efficient and fast experiment (scanning speed and integration time variation, height variation in fast mode), smoothness of steps in experiments and precautions we took.\\
    	%Visualisation: No visualisation needed. we can refer to the intro figure\\
        We have assembled a highly complex machinery for single-molecule detection, but detection does not end at the detectors, and nor does noise. The signal detected by the detectors was digitised using COMPACT RIO and transferred to our computer for readout. To handle the various detection modes, connect detectors and amplifiers, and support sample position and live readout, we developed a custom LabVIEW application. As the refinement of the readout side was not part of this project, I will not elaborate further on the tools used. We used a piezo-controlled 3D-stage (PI), which is low-noise but is not drift-free, so regular repositioning is necessary. During experiments, we opted for a compromise between accurate imaging and time. The most accurate detection would require the sample to move to the laser focus stop, allow detection for the intended integration time, and then move to the next pixel. This makes image acquisition very slow, so we devised a fast-scanning mode in which the sample moves at a constant speed, keeping a single pixel in focus throughout the integration time. This saves speeding up time between pixels, but the stage must speed up and slow down at the edges of the image and recalibrate its x-, y-, and z-positions for every line. As a result of this, the integration time is ever so slightly higher at the edges, and every line has a slight variation in y-position. The images were corrected for this effect. For published images, the more accurate slow mode was used. \\
        The core of this project is to apply electric fields to proteins to measure their dipole moments. Alas, applying said electric fields is not trivial, even with a custom chip and \gls{pcb}. Some trapping experiments required changing parameters during a trapping event, like changing the applied field voltage or carrier frequency. The signal generator we used offered an option to change parameters as needed, but it would update the full set of parameters, not just the one we changed, so the phase of the applied field was not continuous. Therefore, we had to crack a lower-level access which allowed us to change only one parameter, keeping the rest constant. This is not trivial and beyond the scope of this research. I recommend an experienced programmer dedicated to this work. Now, our applied field is smooth in all respects except the changed parameter. 

        \paragraph{To summarise}
        For a complex setup like ours, custom software is a must. For imaging, there is a balance between accuracy and time. For \gls{DEP}-experiments, special attention needs to be given to the smoothness of the applied field. 
        
    %\subsubsection{other noise sources}
    %\label{sec:other}
    %	There are many more things that can happen during assemblage. The following describes various problems that can arise during alignment and noise reduction.
    %		\begin{figure}
    %          	\centering
    %          	\includegraphics[width=\textwidth]{pictures/trubelshooting/NoiseExamples/noiseExamples_almost}
    %          	\caption{\textbf{Example data of various noise and background sources.} 
    %          		We record time traces with a sampling rate of $0,001~ms$ and generate images using a confocal scanning mode. During this data collection process, various noise sources arise and background has to be supressed. Shown here are various traces and images showcasing how noise and background in our setup is experienced. 
    %          		\textbf{A} mechanical noise: \gls{psd} of eliot stage. It is decoupled from the sample exaggerating its vibrations against the sample vibrations causing noise that overshadows any signal we hope to detect.
    %          		\textbf{B} acoustic noise: \gls{psd} of out scattering signal while playing the note $100 Hz$ next to our setup WHICH SAMPLE (GNR OR ),  
    %          		\textbf{C} electrical noise: a broken cable caused the output of our \gls{pcb} to showcase unexpected signals, 
    %          		\textbf{D} detector background, 
    %          		\textbf{E} laser noise: mode hopping of fibre coupled laser, 
    %          		\textbf{F} noise caused by our code: actual speed of our sample during scanning showcased in a 2D-plot. This changes the actual integration time, 
    %          		\textbf{G} noise caused by our sample: sample drying out during scan}
    %          	\label{fig:noiseExamples}
    %        \end{figure}  
%\FloatBarrier
\subsection*{Discussion}
This chapter discussed the construction of a combined \gls{iscat}- and \fl-microscope with \gls{DEP}-trapping capabilities and the problems that can arise during this process. Various sources were mentioned to mitigate the problems, and personal experience was added to this canon.\\
The focus of this chapter was to identify noise sources and eliminate them. But which are the problems to be most vigilant about? Some noise sources are highly situational, such as faulty cables described in \ref{sec:electrical}. But others are more systematic and therefore require careful treatment. First, mechanical noise and alignment go hand in hand. They require the most upkeep and time investment. Differentiating between background that can be subtracted out and noise that disturbs the signal is key here. Next, laser noise, detector background, and general light and vibrational backgrounds can be limited by selecting high-quality equipment and setting it up to ensure longevity. In a similar vein, protection from acoustic noise is most effective when considered during the planning of the laboratory space. The program used to detect both signal and noise serves as the interface to various machines from international producers, and bringing all these aspects together is not trivial. It should be managed by an expert who also regularly updates it. Finally, the sample requires careful handling. In particular, the control of glass thickness cannot be overstated, as the photon budget is limited.\\
In the end, this work helps other experimentalists become more aware of issues and provides guidance on how to address them. 

\subsection*{Outlook}
Several steps to improve and troubleshoot a setup have already been recommended in this chapter, but a setup is never finished and can always be improved. The setup, as presented here, can be improved in the future as well. \\
One of the interceptions necessary would be active drift control. The electric field is highly localised, and even with the PI-stage used to collect this data, drifting away from this spot was observed regularly. Therefore, active drift control should be implemented \cite{niederauerK2}\cite{2015CRISP}\cite{bellve2014Design}. \\
Wide-field \gls{iscat} could be implemented to allow the simultaneous observation of multiple traps. This would increase data output and allow comparisons of datasets under the same conditions. \\
Even with the improvements shown here, a more stable objective and sample mount improve the setup's vibrational stability. An immobilised objective and a lockable sample mount could improve this. This might also help mitigate drift. \\
The background of the gaps could be more efficient. We know that their \scat signal is anisotropic, so placing an analyser in the detection pathway to filter out light from the gaps could help with background suppression.


\newpage
\bibliographystyle{habbrv}
% \bibliographystyle{alpha}
\bibliography{references}
\newpage
\subsection*{Appendix}
% !TeX root = ../main.tex
\label{app:trouble}
\subsubsection{Setup Parts}
	Here the full list off all optical components used in the setup described here.
	\begin{itemize}
		\item [laser] ThorLabs
		\item[14 mirrors] silver mirrors form ThorLabs product number 
	\end{itemize}
\subsubsection{alignment troubleshooting guide}
Here is my adaptation of \hyperref{jonkman2020tutorial}{Jonkman's troubleshooting guide} for imaging and time traces, but do see theirs for more details\cite{jonkman2020tutorial}.
\begin{description}
	\item \textbf{Is it the microscope?}
	\begin{enumerate}
		\item If there is no signal at all, in either the scattering or in fluorescence channels:
		\begin{description}
			\item[Working from back to front]
			\item[Is detection possible?] Are all of the hardware components (main switch/laser, including key, power supplies for detectors and stage) turned on?
			\item[Is the software connected to the detectors and the pi stage?] If not turn them off and on again and run the compact rio connection software
			\item[Is the sample illuminated?] Can you see whether laser light is emitted from the objective lens? This is also a good moment to check if you are in the right focal plane using the bright field mode.
			\item[Is something blocking teh beam path?] Trace the light path from the light source to the sample, and from the sample to the detector, checking each component.
			\item[Are you using enough laser power?] Our setup can anomalously attenuate the laser power using polarisation optics. The power could be set too low to give a signal.
		\end{description}
		\item If the sample appears weak and/or blurry in the z-direction, or you cannot focus:
		\begin{description}
			\item[Is there a bubble in your oil?] Bubble can ocuure when changeing focus rapidly over a bif distance or during every attachemnt of the oil and the sample. If they occure, blowing them away or chaning the oil is enough.
			\item[How does the coma look?] Test this in wide field or on a block placed in the detection path and moving the objective up and down. maybe the objective or sample are mouted tilted. 
			\item[Is the objective hitting the edge?] If you are too close to the edge of a sample the edge of the objective might be stuck on the edge of the sample and not be able to physically move closer.
			\item[Is there something placed in the uncollumated beam path? ] Attenuators and filters placed in an uncollumated beam can displace it. An uncentered beam will cause focal wobble.
			\item[Is the objective clean?] We use lens cleaning tissues and methanol to clean the lens. The thorlabs website has a good viddeo explaining ow to do this\cite{Thorlabs2025OpticsHandling}. this has to be done after every experiment, to avoid damage to the lens surface coating. 
			\item[Is the objective stuck?] Our objective has a spring-loaded front lens, it could become stuck in the compressed position. 
			\item[Could the objective be damaged (oil inside, scratched)? ] Take it out (or ask facility staff to have a look) and check it under a stereo microscope.
		\end{description}
	\end{enumerate}
	\item \textbf{Is it your sample?}
	\begin{enumerate}
		\item Basic checks:
		\begin{description}
			\item[Is the cip the right thickness?] Our objective requires chips and coverslips of 170-$\mu$m-thickness (\# 1.5)
			\item[test another sample] Standard samples for us are electrodes or \gls{gnr}. Can you get a good image now? Are the intensities as expected?
			\item[Has your sample dilution deterioreated?] We have noticed our dilutions deteriourating over time. We suspect the particles clumping and falling out. making a fresh dilution everytime for partices is adviseable.
		\end{description}
		\item Has the sample been damaged or degraded?
		\begin{description}
			\item[Is the galss or the well corrupted?] Especially with our \gls{gnr} samples, which are imaged in HCl this can be dangerous for the objective and need to be detected early. PDMS wells can be leaky and suck particles out. 
			\item[How long ave you imaged?] Our fluorecent particles can bleach over time in the laser and room light. 
		\end{description}
	\end{enumerate}
\end{description}
\FloatBarrier